\vspace{2em}

{
 \normalsize
  \begin{tcolorbox}[teaterminalstyle, title=TEA Program: TOS v1.1.0 Complete Source-Code, breakable]
  \begin{lstlisting}[language=TEA,breaklines=true,numbers=left]
I!:{oooooooo_________________
___oo_____ooooo___ooooo__
___oo____oo____o_oo___oo_
___oo____ooooooo_oo___oo_
___oo____oo______oo___oo_
___oo_____ooooo___oooo_o_
_________________________}                                     
x!:{
TEA OS v1.1.0 is Ready...}
v:vBanner 
L:lPREAMBLE_START
Y:vBanner 
V@:dvLOG_start
V@:dvABOUT_TOS:{ 
#---[ABOUT TEA OS]

NAME: TEA Operating System 
(also "TEA OS" or just "TOS")

ARCHITECT: J. Willrich Lutalo <jwl@nuchwezi.com>

SINCE: 9th APRIL, 2026

INTENT: To manifest a true general purpose, lightweight, minimal but powerful operating system that can essentially just fit inside a single string; to see what the limits of the mathematics of TRANSFORMATICS might be, and what we surely could accomplish with its proper use via modern portable computers.

ABOUT: TEA OS is implemented using the Transforming Executable Alphabet (TEA) programming language. It is based on a specification first laid out in the TEA TAZ. It is a project originating from UGANDA.

PROJECT HOME: Nuchwezi Research | https://tea.nuchwezi.com}
V@:dvOSName:{TEA OS}
V@:dvOSName_short:{TOS}
V@:dvOSVersion:{1.1.0}
V:vSLOGAN:{TEA OS! T can fit in a string!}
I!:{___________________}
V:vDELIM
V@:dvDELIM
I!:{
=o==o==o==o==o==o==o=
}
V:vDELIM2
V@:dvDELIM2
V:vUSERNAME:{user} 
V@:dvLastPS:{}
V@:dvLastCommand:{}
V@:dvLastOutput:{}
Y@:dvFileList
I:{.}
V@:dvFileList
Y@:dvProgramsList
I:{run}
V@:dvProgramsList
Y@:dvUIMODE_MINI
I:{1}
V@:dvUIMODE_MINI
V:fPROMPT_DEFAULT:"
# FIRST: IMPORTANTLY...
v:vMSG #stores AI at invocation
y@:dvDELIM|v:vD2
y:vMSG|x!:{ }|v:vMSG #affix space at prompt end
i!:{} | g*:{
}:vD2:vMSG
v:vMSG
Y@:dfLogAI|v:f # a stealth import :)
y:vMSG|e*:f
#then show log+latest prompt
Y@:dvLOG
i: #returns input..
"
Y:fPROMPT_DEFAULT
V@:dfPROMPT_DEFAULT
V:fPROMPT:"
# FIRST: IMPORTANTLY...
v:vMSG #stores AI at invocation

#check if we should instead be using default ui
y@:dvUIMODE_MINI | f!:^0$:lN_UIMODE_MINI
y@:dfPROMPT_DEFAULT|v:fPROMPT_DEFAULT
y:vMSG | E*:fPROMPT_DEFAULT
q!: # no need to proceed..
l:lN_UIMODE_MINI

#THEN imports...
y@:dfLastIO|v:fLastIO
y@:dfMakePS|v:fMakePS

y@:dvDELIM|v:vD2
y:vMSG|x!:{ }|v:vMSG #affix space at prompt end
i!:{} | g*:{
}:vD2:vMSG
v:vMSG
Y@:dfLogAI|v:f # a stealth import :)
y:vMSG|e*:f
#then show log+latest prompt
#Y@:dvLOG
E*:fMakePS|v:vNEXTPS
E*:fLastIO #fetch last I/O
x!:{
> }|#x*!:vNEXTPS
i: #returns input..
"
Y:fPROMPT
V@:dfPROMPT
V:fLogin:"
#import some functions from database...
y@:dfPROMPT_DEFAULT|v:fPROMPT
y@:dfLogAI|v:fLogAI
y@:dfLogO|v:fLogO
y@:dfGreetUser|v:fGreetUser

# show OS banner
y@:dvLOG_start | v:vSCREEN
e*:fLogAI

#greet user via stored name..
y!@:dvUSERNAME | f!:^0$:lGREET_USER
#Username not yet set, prompt and store
l:lSET_USERNAME
I!:{
Please set a Username:}|v:vPromptMSG
e*:fLogAI
y:vSCREEN | x*!:vPromptMSG | i:
t!.:|v:|v!:|f:^0$:lSET_USERNAME
y:|V@:dvUSERNAME # store in database
e*:fLogAI

l:lGREET_USER
e*:fGreetUser|v:vGreetings
e*:fLogO | e*:fPROMPT
"
V:fGreetUser:"
y@:dvUSERNAME
x:{
Hello, }|x!:{. Welcome to TOS!}
"
Y:fGreetUser
V@:dfGreetUser
V:fOSHandle:"
y@:dvOSName | v:vN
y@:dvOSVersion | v:vV
g*:{ v}:vN:vV
"
Y:fOSHandle
V@:dfOSHandle
V:fMakePS:"
z.:TIME |v:vT
y@:dvOSName_short | v:vPWD
v:vD:{>}
g*:{/}:vPWD:vT:vD
"
Y:fMakePS
V@:dfMakePS
V:fLogIO:"
#load+execute function from database! 
y@:dfLastIO|e:
v:vLastIO
y@:dvLOG | v:vLOG
g*:{
}:vLOG:vLastIO | v@:dvLOG
"
V:fLogAI:"
v:vAI
y@:dvLOG | v:vLOG
g*:{
}:vLOG:vAI | v@:dvLOG
#return AI
y:vAI
"
Y:fLogAI
v@:dfLogAI
V:fLogO:"
v@:dvLastOutput
v@:dvDATABUFFER #also stash into special data buffer
Y@:dvLastOutput
"
V:fLastIO:"
y@:dvLastPS | v:vPS
y@:dvLastCommand | v:vC
y@:dvLastOutput | v:vO
g*:{}:vPS:vC|v:vCC
g*:{
}:vCC:vO
"
y:fLastIO
v@:dfLastIO
L:lPREAMBLE_END
L:lSYS_UTILITY_START
V:fFixURL:"
v:vURL|z:
f:{^https?://}:lFINE
y:vURL|x:{https://}
v:vURL
l:lFINE
y:vURL
"
V:fResetLog:"
y@:dvLastOutput | v@:dvLastOutput_backup
c!@:dvLOG
c*:vLOG
C!@:dvLastOutput
y@:dfGreetUser|E:|v:vGreetings
y@:dvLOG_start
x*!:vGreetings
"
Y:fResetLog
V@:dfResetLog
V:fGetLatestOutputBuffer:"
y@:dvLastOutput_backup
f!:^$:lN_USED_BACKUP
y@:dvLastOutput
l:lN_USED_BACKUP
c@:dvLastOutput_backup
"
Y:fGetLatestOutputBuffer
V@:dfGetLatestOutputBuffer
L:lSYS_UTILITY_END
L:lSYS_HELP_START
V:fGetUIMenu:"
y@:dfOSHandle|e:|x!:{ INTERFACE System Commands:
uimini
uilog}
v:vMSG
y@:dvDELIM2|v:vD
g*:{}:vD:vMSG:vD
"
V:fGetFileCMDMenu:"
y@:dfOSHandle|e:|x!:{ FILE System Commands:
create FILENAME
init FILENAME
write FILENAME
(rename | rn | mv | move) FILENAME NEWNAME
(copy|clone) FILE1 FILE2
exists FILENAME
find PATTERN
(show | cat | read) FILENAME
export FILENAME
(del | delete | rm) FILENAME
(ls | list | dir | files) PATTERN}
v:vMSG
y@:dvDELIM2|v:vD
g*:{}:vD:vMSG:vD
"
V:fGetProgrammingMenu:"
y@:dfOSHandle|e:|x!:{ PROGRAMMING System Commands:
(runnable | register | add-program) FILENAME
(isrunnable | canrun) FILENAME
(run | exec | execute) FILENAME
(del-prog | delete-program | unregister) FILENAME
(update | update-program) FILENAME
PROGRAM
(help | man | whatis | info | about) PROGRAM}
v:vMSG
y@:dvDELIM2|v:vD
g*:{}:vD:vMSG:vD
"
V:fGetUtilityMenu:"
y@:dfOSHandle|e:|x!:{ UTILITY System Commands:
clear | reset
date
echo MESSAGE
exit
export
(fetch | web) [URL] [FILENAME]
logout
math EXPRESSION
print
tea CODE
time
timestamp
username NAME
whoami | name}
v:vMSG
y@:dvDELIM2|v:vD
g*:{}:vD:vMSG:vD
"
V:fGetHelpMenu:"
y@:dfOSHandle|e:|x!:{ HELP System Commands:
about
(help | man | whatis | info | about) [NAME]
help-file
help-help
help-prog
help-ui
help-util}
v:vMSG
y@:dvDELIM2|v:vD
g*:{}:vD:vMSG:vD
"
V:fGetAboutTOS:"
y@:dfOSHandle|e:|x:{ABOUT: }
v:vMSG
y@:dvABOUT_TOS|x*:vMSG|v:vMSG
y@:dvDELIM2|v:vD
g*:{}:vD:vMSG:vD
"
V:fGetHelpOn:"
#asumes topic is passed via AI
v:vHelpTOPIC

#some imports...
y@:dfGetFilenameProtocolMSG|v:fGetFilenameProtocolMSG
E*:fGetFilenameProtocolMSG | v:vFILENAME_PROTOCOL

#set default..
r@:{That help topic [#vHelpTOPIC#] is not understood or supported. Try the 'help' command to see what help commands there are. For example 'help whatis'}|v:vMSG

y:vHelpTOPIC

#> help about
f!:^about$:lN_h0
v:vMSG:{about <~ Displays basic info about the currently running version of the TEA Operating System.}
j:lh1_present
l:lN_h0

#> help whatis
f!:^whatis$:lN_h1
v:vMSG:{whatis NAME <~ Displays help info about a given system command or user program.}
j:lh1_present
l:lN_h1

#> help man
f!:^man$:lN_h1b
v:vMSG:{man NAME <~ Displays help info about a given system command or user program. It is just another alias for 'whatis' or 'help'}
j:lh1_present
l:lN_h1b

#> help help
f!:^help$:lN_h2
v:vMSG:{Present menu for user to access help about all available [essential] commands based on category.}
j:lh1_present
l:lN_h2

f!:^help-file$:lN_h3
v:vMSG:{Display all available commands in the FILE I/O category or sub-system.}
j:lh1_present
l:lN_h3

f!:^help-prog$:lN_h4
v:vMSG:{Display all available commands in the PROGRAMMING category or sub-system.}
j:lh1_present
l:lN_h4

f!:^help-util$:lN_h5
v:vMSG:{Display all available commands in the UTILITIES category or sub-system.}
j:lh1_present
l:lN_h5

f!:^help-help$:lN_h6
v:vMSG:{Display help on how to use the HELP facility or sub-system.}
j:lh1_present
l:lN_h6

f!:{^(clear|reset)$}:lN_hu0
v:vCMD
r@:{#vCMD#  <~ Resets and clears the current TEA OS session log since login.} | v:vMSG
j:lh1_present
l:lN_hu0

f!:{^(fetch|web)$}:lN_hu1
v:vCMD
r@:{#vCMD# URL <~ fetch and output contents of the specified [online/network] resource.

#vCMD# URL FILENAME <~ fetch and save contents at URL into file record named FILENAME. #vFILENAME_PROTOCOL#} | v:vMSG
j:lh1_present
l:lN_hu1

f!:^tea$:lN_hu2
v:vMSG:{tea CODE.. <~ take everything following the keyword 'tea' as a TEA program, run and output the result: It's a simple way to write and run basic TEA programs (TEA-one-liners) in TOS.}
j:lh1_present
l:lN_hu2

f!:^math$:lN_hu3
v:vMSG:{math EXPRESSION <~ execute basic mathematics on the command-line as specified in EXPRESSION.}
j:lh1_present
l:lN_hu3

f!:^date$:lN_hu4
v:vMSG:{date <~ Show the current date.}
j:lh1_present
l:lN_hu4

f!:^time$:lN_hu5
v:vMSG:{time <~ Show the current time.}
j:lh1_present
l:lN_hu5

f!:^timestamp$:lN_hu6
v:vMSG:{timestamp <~ Show the current timestamp.}
j:lh1_present
l:lN_hu6

f!:^print$:lN_hu7
v:vMSG:{print <~ Show or preview current output buffer contents (usually, output of last command). After clearing log, but also after exit/system shutdown/login also shows final output before the log was reset.}
j:lh1_present
l:lN_hu7

f!:^export$:lN_hu8
r@:{export <~ Take the output of the last command or rather, current contents of the output buffer and export that outside the operating system (also terminates the operating system session).

export FILENAME <~ Exit the system, outputting the contents of given file. #vFILENAME_PROTOCOL#} | v:vMSG
j:lh1_present
l:lN_hu8

f!:^exit$:lN_hu9
v:vMSG:{exit <~ Ends the TOS session returning nothing.}
j:lh1_present
l:lN_hu9

f!:^logout$:lN_hu10
v:vMSG:{logout <~ Resets username and system logs and ends the TOS session returning nothing.}
j:lh1_present
l:lN_hu10

f!:^username$:lN_hu11
v:vMSG:{username NAME <~ Overrides current username with specified value.}
j:lh1_present
l:lN_hu11

f!:^help-ui$:lN_h12
v:vMSG:{Display all available commands in the USER INTERFACE category or sub-system.}
j:lh1_present
l:lN_h12

f!:^uimini$:lN_h13
v:vMSG:{uimini <~ Configure TEA OS interface to render using mini-terminal mode. Setting persists across logins until changed. This mode is suitable for environments such as using TOS on some Desktop Computer Systems.}
j:lh1_present
l:lN_h13

f!:^uilog$:lN_h14
v:vMSG:{uilog <~ Configure TEA OS interface to render using long/log-terminal mode. Setting persists across logins until changed. This mode is suitable for environments such as using TOS on the TEA MOBILE IDE.}
j:lh1_present
l:lN_h14

f!:^whoami$:lN_h15
v:vCMD
r@:{#vCMD# <~ Display the NAME of currently active USER.} | v:vMSG
j:lh1_present
l:lN_h15

f!:^create$:lN_h16
r@:{create FILENAME <~ Creates blank file with name FILENAME. #vFILENAME_PROTOCOL#}|v:vMSG
j:lh1_present
l:lN_h16

f!:{^(ls|list|dir|files)$}:lN_h17
v:vCMD
r@:{#vCMD# <~ Display list of ALL available files.

#vCMD# PATTERN <~ Display list of ONLY files matching regular expression PATTERN.}
| v:vMSG
j:lh1_present
l:lN_h17

f!:^init$:lN_h18
r@:{init FILENAME <~ Create a file using the given name and the currently available input/last-output (or the PRINT buffer) or blank if none. #vFILENAME_PROTOCOL#}|v:vMSG
j:lh1_present
l:lN_h18

f!:^write$:lN_h19
r@:{write FILENAME <~ Write the contents of the current available input/last-output (or the PRINT buffer) or blank if none, to a file with the specified name. If the file already exists, essentially overwrites its contents. #vFILENAME_PROTOCOL#} | v:vMSG
j:lh1_present
l:lN_h19

f!:{^(show|cat|read)$}:lN_h20
v:vCMD
r@:{#vCMD# FILENAME <~ Preview the contents of given file. #vFILENAME_PROTOCOL#} | v:vMSG
j:lh1_present
l:lN_h20

f!:^echo$:lN_h21
v:vMSG:{echo MESSAGE <~ Simply print MESSAGE as is.}
j:lh1_present
l:lN_h21

f!:{^(find)$}:lN_h22
v:vCMD
r@:{#vCMD# PATTERN <~ Display list of ONLY files matching regular expression PATTERN.}
| v:vMSG
j:lh1_present
l:lN_h22

f!:{^(del|delete|rm)$}:lN_h23
v:vCMD
r@:{#vCMD# PATTERN <~ Remove from storage, any and all files whose name matches PATTERN. CAN NOT DELETE registered user program files though.} | v:vMSG
j:lh1_present
l:lN_h23

f!:{^(rename|rn|mv|move)$}:lN_h24
v:vCMD
r@:{#vCMD# FILENAME NEWNAME <~ Change the name of an existing file from FILENAME to NEWNAME. WILL ALSO update the registered [user] programs list if FILENAME was a registered program. #vFILENAME_PROTOCOL#} | v:vMSG
j:lh1_present
l:lN_h24

f!:{^(exists)$}:lN_h25
v:vCMD
r@:{#vCMD# FILENAME <~ Show whether or not a file named FILENAME exists. #vFILENAME_PROTOCOL#} | v:vMSG
j:lh1_present
l:lN_h25

f!:{^(copy|clone)$}:lN_h26
v:vCMD
r@:{#vCMD# FILENAME FILE2 <~ Copies contents of FILENAME to FILE2 creating the destination if it doesn't exist. #vFILENAME_PROTOCOL#} | v:vMSG
j:lh1_present
l:lN_h26

f!:{^(runnable|register|add-program)$}:lN_h27
v:vCMD
r@:{#vCMD# FILENAME <~ Recognize FILENAME object as an executable [user] program or command. FILENAME must already be an existing TEA file. #vFILENAME_PROTOCOL#} | v:vMSG
j:lh1_present
l:lN_h27

f!:{^programs$}:lN_h28
v:vCMD
r@:{#vCMD# <~ Show list of all available registered [TEA]  user programs by their registered names.

#vCMD# PATTERN <~ With PATTERN specified, ONLY show those programs whose names match regular expression PATTERN.}
| v:vMSG
j:lh1_present
l:lN_h28

f!:{^(run|exec|execute)$}:lN_h29
v:vCMD
r@:{#vCMD# FILENAME <~ Execute the given file record as a [TEA] program. #vFILENAME_PROTOCOL#} | v:vMSG
j:lh1_present
l:lN_h29

f!:{^(isrunnable|canrun)$}:lN_h30
v:vCMD
r@:{#vCMD# FILENAME <~ Show whether or not a file named FILENAME can be executed as a [TEA] program. #vFILENAME_PROTOCOL#} | v:vMSG
j:lh1_present
l:lN_h30

f!:{^(del-prog|delete-program|unregister)$}:lN_h31
v:vCMD
r@:{#vCMD# FILENAME <~ Revokes status of FILENAME as an executable user program (doesn't actually delete the underlying file object though). #vFILENAME_PROTOCOL#} | v:vMSG
j:lh1_present
l:lN_h31


f!:{^(update|update-program)$}:lN_h32
v:vCMD
r@:{#vCMD# PROGRAMFILE FILENAME <~ Replace/override original program named PROGRAMFILE with contents of FILENAME (updates a program's code [in PROGRAMFILE] preserving its existing name). Both files PROGRAMFILE and FILENAME must actually exist for this to work, but only PROGRAMFILE needs to be a registered program. #vFILENAME_PROTOCOL#} | v:vMSG
j:lh1_present
l:lN_h32

f!:{^(program)$}:lN_h33
v:vCMD | z!: | v:vCMD #use upercase..
r@:{#vCMD# <~ Executes #vCMD# as a system command if it exists --- first, if it is any of the standard commands in TEA OS, otherwise if it is a valid filename "#vCMD#" that is already recognized/registered as a user-created program - this later case similar to calling "run #vCMD#", otherwise shows error message command-not-found.} | v:vMSG
j:lh1_present
l:lN_h33

# the user-program catch-all...
f!:{^([a-z0-9_]+(\.[a-z0-9]+)?)$}:lN_h34
v:vPROGRAMFILE 
y@:dfProgramExists | v:fProgramExists #import
y:vPROGRAMFILE | e*:fProgramExists 
f!:{^PROGRAM EXISTS$}:lN_h34
r@:{#vPROGRAMFILE# <~ Is a valid user-registered [TEA] program. You might run or execute it by directly invoking "#vPROGRAMFILE#" or via "run #vPROGRAMFILE#".} | v:vMSG
j:lh1_present
l:lN_h34

l:lh1_present
y@:dvDELIM2|v:vD
g*:{}:vD:vMSG:vD
"
L:lSYS_HELP_END
L:lSYS_FILE_START
V:fGetFilenameProtocolMSG:"i!:{Valid FILENAME should consist of ONLY letters a-zA-Z, numbers 0-9, underscore _ and/or a dot .}"
Y:fGetFilenameProtocolMSG
V@:dfGetFilenameProtocolMSG
V:fGenDBFilename:"x:{DFS_}"
Y:fGenDBFilename
V@:dfGenDBFilename
V:fDecodeDBFilename:"d:^DFS_"
Y:fDecodeDBFilename
V@:dfDecodeDBFilename
V:fIsFileNameValid:"
f!:^[a-z0-9_]+(\.[a-z0-9]+)?$:lBAD_FILE_NAME
i!:1|q!:
l:lBAD_FILE_NAME
i!:0
"
Y:fIsFileNameValid
V@:dfIsFileNameValid
V:fCreateFile:"
# first, grab incoming function argument...
v:vFILENAME|t!.:|v:vFNAME

#then import functions..
y@:dfGenDBFilename|v:fGenDBFilename
y@:dfGetLatestOutputBuffer|v:fGetLatestOutputBuffer
y@:dfIsFileNameValid|v:fIsFileNameValid

y:vFNAME | f:^$:lNFINE # filename shouldn't be blank
# file name should conform to accepted patterns
y:vFNAME | E*:fIsFileNameValid | f!:0:lGOOD_FILENAME
r@:{Invalid File Name [#vFNAME#]!
FILE NOT CREATED.}|q!:
l:lGOOD_FILENAME

# fetch current file name list..
y@:dvFileList | v:vFileList
# ensure file name doesn't yet exist in records
#e.g v:vFileList:{F1|F2|F33|File|File}
#v:vFNAME:{f1}
# check if file name is unique
r@:{(^#vFNAME#\|)|(\|#vFNAME#\|)|(\|#vFNAME#$)}|v:vFNAME_REGEX
y:vFileList | f*!:vFNAME_REGEX:lFNAME_UNIQUE:lFNAME_DUPLICATE
l:lFNAME_UNIQUE
# add name to file list 
g*:{|}:vFileList:vFNAME | v:vFileList | v@:dvFileList 
#also create blank file record in db filesystem
y:vFNAME|E*:fGenDBFilename|v:vdFNAME

v@:dvDATA:{} # we'll write empty file
r@:{y@:dvDATA|v@:#vdFNAME#}|E:

i!:{CREATED NEW BLANK FILE: }|x*!:vFNAME
q!:
l:lFNAME_DUPLICATE
i!:{File Already Exists! 
NOT DUPLICATING FILE: }|x*!:vFNAME
q!:
l:lNFINE
i!:{Invalid [blank] File Name!
FILE NOT CREATED.}
"
Y:fCreateFile
V@:dfCreateFile
V:fDeleteFile:"
# first, grab incoming function argument...
v:vFILENAME|t!.:|v:vFNAME

#then import functions..
y@:dfGenDBFilename|v:fGenDBFilename
y@:dfIsFileNameValid|v:fIsFileNameValid

y:vFNAME | f:^$:lNFINE # filename shouldn't be blank
# file name should conform to accepted patterns
y:vFNAME | E*:fIsFileNameValid | f!:0:lGOOD_FILENAME
r@:{Invalid File Name [#vFNAME#]!
FILE NOT DELETED.}|q!:
l:lGOOD_FILENAME

# fetch current file name list..
y@:dvFileList | v:vFileList
# remove file name from records

#v:vFileList:{F1|F2|F33|F3|File|File}
#v:vFNAME:{File}
# get filename pattern

# create delete filename patterns + apply them..
r@:{(^#vFNAME#\|)|(\|#vFNAME#$)}|v:vFNAME_REGEX_ENDS
r@:{(\|#vFNAME#\|)}|v:vFNAME_REGEX_MID
r@:{^#vFNAME#$}|v:vFNAME_REGEX_ALL
# use it to del
y:vFileList | d*:vFNAME_REGEX_ENDS | v:vFileList
y:vFileList | d*:vFNAME_REGEX_ALL | v:vFileList
v:vDIV:{|}
r*:vFileList:vFNAME_REGEX_MID:vDIV | v:vFileList

# update file list 
y:vFileList | v@:dvFileList 
#also delete file record from db filesystem
y:vFNAME|E*:fGenDBFilename|v:vdFNAME
r@:{c!@:#vdFNAME#}|E:

i!:{DELETED FILE: }|x*!:vFNAME
q!:

l:lNFINE
i!:{Invalid [blank] File Name!
FILE NOT DELETED.}
"
Y:fDeleteFile
V@:dfDeleteFile
V:fShowFilesList:"
# takes no parameters ;)
# fetch current file name list..
y@:dvFileList
#process and present..
h!:{\|}|d:{\|}
"
V:fSearchFilesList:"
# takes file regex pattern parameter
v:vREGEX_PATTERN
# fetch current file name list..
y@:dvFileList
#process and present..
h!:{\|}|d:{\|}
# only keep what matches...
k*:vREGEX_PATTERN
"
V:fDeleteFilesList:"
# takes file regex pattern parameter
v:vREGEX_PATTERN

#imports..
y@:dfDeleteFile | v:fDeleteFile
y@:dfProgramExists | v:fProgramExists

# fetch current file name list..
y@:dvFileList
#process and present..
h!:{\|}|d:{\|}
# only keep what matches...
k*:vREGEX_PATTERN 
g:{|} | v:vDelFileList | v:vDelFileListORIG
f:^$:NO_FILES_TO_DELETE
# to track files we couldn't delete
v:vNOTDelFileList:{}

# go through files and delete..
v:vNCOUNTER:0
v:vNCOUNTER_DEL:0

l:lSTART_DEL
y:vDelFileList
f:^$:lDONE_DEL
#fetch..
y:vDelFileList | d!:{^[^\|]+\|?} 
d:{\|} | v:vDelFileNAME 

#if actually blank.. go pop
f:^$:lPOPDEL_LIST

# ensure it is not a registered program
y:vDelFileNAME | e*:fProgramExists | f:{^PROGRAM EXISTS$}:lDELETE_FAILED

#r@:{--Continuing to delete file [#vDelFileNAME#]} | q!:

# delete file..
y:vDelFileNAME | e*:fDeleteFile
#q!: # check...

f!:{DELETED FILE}:lDELETE_FAILED:lDELETED_OK
l:lDELETE_FAILED
# add name to NOT deleted file list 
g*:{|}:vNOTDelFileList:vDelFileNAME  | v:vNOTDelFileList
j:lDONE_DELETED_OK
l:lDELETED_OK
y:vNCOUNTER_DEL | x!:+1 | r.: | v:vNCOUNTER_DEL
l:lDONE_DELETED_OK
y:vNCOUNTER | x!:+1 | r.: | v:vNCOUNTER

#pop..
l:lPOPDEL_LIST
y:vDelFileList | d:{^[^\|]+\|?} | v:vDelFileList
#x:{Remaining: } | q!:
j:lSTART_DEL # iterate

l:lDONE_DEL
y:vDelFileListORIG|d:{^\|}:{\|$}|v:vDelFileListORIG
y:vNOTDelFileList|d:{^\|}:{\|$}|v:vNOTDelFileList
r@:{#vNCOUNTER# - #vNCOUNTER_DEL#}|r.:|v:vNCOUNTER_NDEL
y:vNCOUNTER_DEL | f:^1$:lDEL_REP_1
r@:{DELETED #vNCOUNTER_DEL# file[s] from list:
#vDelFileListORIG#
EXCEPT #vNCOUNTER_NDEL#:
#vNOTDelFileList#}|q!:

l:lDEL_REP_1
y:vDelFileListORIG|d:{^\|}:{\|$}|v:vDelFileListORIG
y:vNCOUNTER | f!:^1$:lDEL_REP_2
r@:{DELETED ONLY #vNCOUNTER_DEL# file:
#vDelFileListORIG#
}|q!:

l:lDEL_REP_2
y:vDelFileListORIG|d:{^\|}:{\|$}|v:vDelFileListORIG
y:vNOTDelFileList|d:{^\|}:{\|$}|v:vNOTDelFileList
r@:{DELETED ONLY #vNCOUNTER_DEL# file from list:
#vDelFileListORIG#
EXCEPT #vNCOUNTER_NDEL#:
#vNOTDelFileList#}|q!:

l:NO_FILES_TO_DELETE
r@:{#vREGEX_PATTERN# matched no FILES!
NO FILE WAS DELETED}
"
V:fInitFile:"
v:vFILENAME|t!.:|v:vFNAME

#imports...
y@:dfCreateFile|v:fCreateFile
y@:dfGenDBFilename|v:fGenDBFilename
y@:dfGetLatestOutputBuffer|v:fGetLatestOutputBuffer
y@:dfIsFileNameValid|v:fIsFileNameValid

y:vFNAME | E*:fIsFileNameValid | f!:0:lGOOD_FILENAME
r@:{Invalid File Name [#vFNAME#]!
FILE NOT CREATED.}|q!:
l:lGOOD_FILENAME

# then, try to create the file...
y:vFNAME | E*:fCreateFile
# check if it was created...
f!:{CREATED NEW BLANK FILE}:lFILE_NOT_CREATED
# get name of that newly created file...
y:vFNAME|E*:fGenDBFilename|v:vdFNAME

# get current output buffer
y@:dvDATABUFFER | v@:dvDATA

# only if it was empty then...
f!:^$:lFINE_DATA_BUFFER
# use print buffer...
e*:fGetLatestOutputBuffer | v@:dvDATA
l:lFINE_DATA_BUFFER

# store it into file system record...
r@:{y@:dvDATA|v@:#vdFNAME#}|E:

#a test...
#r@:{y@:#vdFNAME#}|E:
#x:{File now contains:
#}|q!:

i!:{CREATED and INITIALIZED NEW File: }|x*!:vFNAME
q!:
l:lFILE_NOT_CREATED
i!:{FAILED TO CREATE FILE: } | x*!:vFNAME
"
V:fWriteFile:"
v:vFILENAME|t!.:|v:vFNAME

#imports...
y@:dfCreateFile|v:fCreateFile
y@:dfGenDBFilename|v:fGenDBFilename
y@:dfGetLatestOutputBuffer|v:fGetLatestOutputBuffer
y@:dfIsFileNameValid|v:fIsFileNameValid

y:vFNAME | E*:fIsFileNameValid | f!:0:lGOOD_FILENAME
r@:{Invalid File Name [#vFNAME#]!
FILE NOT WRITTEN.}|q!:
l:lGOOD_FILENAME

# then, try to create the file...
y:vFNAME | E*:fCreateFile
# check if it was created...
f!:{(CREATED NEW BLANK FILE)|(NOT DUPLICATING FILE)}:lFILE_NOT_CREATED
# get name of that file...
y:vFNAME|E*:fGenDBFilename|v:vdFNAME
# get current output buffer
y@:dvDATABUFFER | v@:dvDATA

# only if it was empty then...
f!:^$:lFINE_DATA_BUFFER
# use print buffer...
e*:fGetLatestOutputBuffer | v@:dvDATA
l:lFINE_DATA_BUFFER

# store it into file system record...
r@:{y@:dvDATA|v@:#vdFNAME#}|E:

i!:{DATA WRITTEN to File: }|x*!:vFNAME
q!:
l:lFILE_NOT_CREATED
i!:{FAILED TO WRITE FILE: } | x*!:vFNAME
"
V:fSaveDataToFile:"
v:vFILENAME|t!.:|v:vFNAME

#imports...
y@:dfCreateFile|v:fCreateFile
y@:dfGenDBFilename|v:fGenDBFilename
y@:dfGetLatestOutputBuffer|v:fGetLatestOutputBuffer
y@:dfIsFileNameValid|v:fIsFileNameValid

y:vFNAME | E*:fIsFileNameValid | f!:0:lGOOD_FILENAME
r@:{Invalid File Name [#vFNAME#]!
FILE NOT WRITTEN.}|q!:
l:lGOOD_FILENAME

# then, try to create the file...
y:vFNAME | E*:fCreateFile
# check if it was created...
f!:{(CREATED NEW BLANK FILE)|(NOT DUPLICATING FILE)}:lFILE_NOT_CREATED
# get name of that file...
y:vFNAME|E*:fGenDBFilename|v:vdFNAME
# get current data buffer
# and store it into file system record...
r@:{y@:dvDATABUFFER|v@:#vdFNAME#}|E:

i!:{DATA WRITTEN to File: }|x*!:vFNAME
q!:
l:lFILE_NOT_CREATED
i!:{FAILED TO WRITE FILE: } | x*!:vFNAME
"
V:fReadFile:"
v:vFILENAME|t!.:|v:vFNAME

#imports...
y@:dfGenDBFilename|v:fGenDBFilename
y@:dfIsFileNameValid|v:fIsFileNameValid

y:vFNAME | E*:fIsFileNameValid | f!:0:lGOOD_FILENAME
r@:{Invalid File Name [#vFNAME#]!
FILE NOT READ.}|q!:
l:lGOOD_FILENAME

# then, try to read the file...
y:vFNAME|E*:fGenDBFilename|v:vdFNAME

r@:{y@:#vdFNAME#}|E:
"
y:fReadFile 
v@:dfReadFile
V:fFileExists:"
# first, grab incoming function argument...
v:vFILENAME|t!.:|v:vFNAME

#then import functions..
y@:dfGenDBFilename|v:fGenDBFilename
y@:dfIsFileNameValid|v:fIsFileNameValid

y:vFNAME | f:^$:lNFINE
y:vFNAME | E*:fIsFileNameValid | f!:0:lGOOD_FILENAME
r@:{Invalid File Name [#vFNAME#]!
FILE MAY NOT EXIST.}|q!:
l:lGOOD_FILENAME

# fetch current file name list..
y@:dvFileList | v:vFileList


#v:vFileList:{F1|F2|F33|F3|File|File}
#v:vFNAME:{f1}

# create filename patterns..
r@:{(^#vFNAME#\|)|(\|#vFNAME#$)}|v:vFNAME_REGEX_ENDS
r@:{(\|#vFNAME#\|)}|v:vFNAME_REGEX_MID
r@:{^#vFNAME#$}|v:vFNAME_REGEX_ALL
# use them to test for existence
y:vFileList | f*:vFNAME_REGEX_ENDS:lFILE_EXISTS
f*:vFNAME_REGEX_ALL:lFILE_EXISTS
f*:vFNAME_REGEX_MID:lFILE_EXISTS

i!:{FILE DOES NOT EXIST}|q!:

l:lFILE_EXISTS
i!:{FILE EXISTS}
q!:

l:lNFINE
i!:{Invalid [blank] File Name!
FILE DOES NOT EXIST.}
"
L:lSYS_FILE_END
L:lSYS_PROG_START
V:fRegisterProgramFile:"
# we assume given program file already exists!
# first, grab incoming function argument...
v:vFILENAME|t!.:|v:vFNAME

#then import functions..
y@:dfGenDBFilename|v:fGenDBFilename
y@:dfIsFileNameValid|v:fIsFileNameValid

y:vFNAME | f:^$:lNFINE
y:vFNAME | E*:fIsFileNameValid | f!:0:lGOOD_FILENAME
r@:{Invalid File Name [#vFNAME#]!
CANNOT REGISTER PROGRAM.}|q!:
l:lGOOD_FILENAME

# fetch current program file name list..
y@:dvProgramsList | v:vProgramsList

#v:vFileList:{F1|F2|F33|F3|File|File}
#v:vFNAME:{f1}

# check if program name is unique:
# create filename patterns...
r@:{(^#vFNAME#\|)|(\|#vFNAME#\|)|(\|#vFNAME#$)}|v:vFNAME_REGEX
# use them to test for program existence
y:vProgramsList | f*!:vFNAME_REGEX:lPROGNAME_UNIQUE:lPROG_EXISTS
l:lPROGNAME_UNIQUE
# So, program isn't registered yet, meaning its ok
# to proceed to register it...
# add name to programs list 
g*:{|}:vProgramsList:vFNAME | v:vProgramsList | v@:dvProgramsList
r@:{User PROGRAM [#vFNAME#] is now REGISTERED.
CAN BE INVOKED as SYSTEM COMMAND.} | q!:

l:lPROG_EXISTS
r@:{User PROGRAM [#vFNAME#] Already EXISTS!
CANNOT re-REGISTER PROGRAM.} | q!:

l:lNFINE
i!:{Invalid [blank] File Name!
CANNOT REGISTER PROGRAM.}
"
V:fShowProgramsList:"
y@:dvProgramsList|h!:{\|}|d:{\|}
"
V:fSearchProgramsList:"
v:vREGEX_PATTERN
y@:dvProgramsList|h!:{\|}|d:{\|}|k*:vREGEX_PATTERN
"
V:fRunFile:"
v:vFILENAME|t!.:|v:vFNAME

#imports...
y@:dfGenDBFilename|v:fGenDBFilename
y@:dfIsFileNameValid|v:fIsFileNameValid

y:vFNAME | E*:fIsFileNameValid | f!:0:lGOOD_FILENAME
r@:{Invalid File Name [#vFNAME#]!
FILE NOT EXECUTED.}|q!:
l:lGOOD_FILENAME

# then, try to run the file...
y:vFNAME|E*:fGenDBFilename|v:vdFNAME

# for now, run it as a TEA program..
r@:{y@:#vdFNAME#}|E:|E:
"
V:fProgramExists:"
v:vFILENAME|t!.:|v:vFNAME

y@:dfGenDBFilename|v:fGenDBFilename
y@:dfIsFileNameValid|v:fIsFileNameValid

y:vFNAME | f:^$:lNFINE
y:vFNAME | E*:fIsFileNameValid | f!:0:lGOOD_FILENAME
r@:{Invalid File Name [#vFNAME#]!
PROGRAM MAY NOT EXIST.}|q!:
l:lGOOD_FILENAME

# fetch current programs name list..
y@:dvProgramsList | v:vProgramsList

r@:{(^#vFNAME#\|)|(\|#vFNAME#\|)|(\|#vFNAME#$)}|v:vFNAME_REGEX
y:vProgramsList | f*:vFNAME_REGEX:lPROG_EXISTS

i!:{PROGRAM DOES NOT EXIST}|q!:

l:lPROG_EXISTS
i!:{PROGRAM EXISTS}|q!:

l:lNFINE
i!:{Invalid [blank] File Name!
PROGRAM DOES NOT EXIST.}
"
Y:fProgramExists 
V@:dfProgramExists
V:fUnRegisterProgram:"
# we assume given program file already exists!
# first, grab incoming function argument...
v:vFILENAME|t!.:|v:vFNAME

#then import functions..
y@:dfGenDBFilename|v:fGenDBFilename
y@:dfIsFileNameValid|v:fIsFileNameValid

y:vFNAME | f:^$:lNFINE
y:vFNAME | E*:fIsFileNameValid | f!:0:lGOOD_FILENAME
r@:{Invalid File Name [#vFNAME#]!
CANNOT UNREGISTER PROGRAM.}|q!:
l:lGOOD_FILENAME

# fetch current program file name list..
y@:dvProgramsList | v:vProgramsList

# create delete filename patterns + apply them..
r@:{(^#vFNAME#\|)|(\|#vFNAME#$)}|v:vFNAME_REGEX_ENDS
r@:{(\|#vFNAME#\|)}|v:vFNAME_REGEX_MID
r@:{^#vFNAME#$}|v:vFNAME_REGEX_ALL
# use it to del
y:vProgramsList | d*:vFNAME_REGEX_ENDS | v:vProgramsList
y:vProgramsList| d*:vFNAME_REGEX_ALL | v:vProgramsList
v:vDIV:{|}
r*:vProgramsList:vFNAME_REGEX_MID:vDIV | v:vProgramsList

# update programs list 
y:vProgramsList| v@:dvProgramsList 

i!:{UNREGISTERED PROGRAM FILE: }|x*!:vFNAME
q!:

l:lNFINE
i!:{Invalid [blank] File Name!
PROGRAM FILE NOT UNREGISTERED.}
"
Y:fUnRegisterProgram 
V@:dfUnRegisterProgram
V:fDeleteProgramFile:"
# first, grab incoming function argument...
v:vFILENAME|t!.:|v:vFNAME

#then import functions..
y@:dfGenDBFilename|v:fGenDBFilename
y@:dfIsFileNameValid|v:fIsFileNameValid
y@:dfUnRegisterProgram |v:fUnRegisterProgram

y:vFNAME | f:^$:lNFINE # filename shouldn't be blank
# file name should conform to accepted patterns
y:vFNAME | E*:fIsFileNameValid | f!:0:lGOOD_FILENAME
r@:{Invalid File Name [#vFNAME#]!
PROGRAM FILE NOT DELETED.}|q!:
l:lGOOD_FILENAME

# fetch current file name list..
y@:dvFileList | v:vFileList


# create delete filename patterns + apply them..
r@:{(^#vFNAME#\|)|(\|#vFNAME#$)}|v:vFNAME_REGEX_ENDS
r@:{(\|#vFNAME#\|)}|v:vFNAME_REGEX_MID
r@:{^#vFNAME#$}|v:vFNAME_REGEX_ALL
# use it to del
y:vFileList | d*:vFNAME_REGEX_ENDS | v:vFileList
y:vFileList | d*:vFNAME_REGEX_ALL | v:vFileList
v:vDIV:{|}
r*:vFileList:vFNAME_REGEX_MID:vDIV | v:vFileList

# update file list 
y:vFileList | v@:dvFileList 
#also delete file record from db filesystem
y:vFNAME|e*:fGenDBFilename|v:vdFNAME
r@:{c!@:#vdFNAME#}|E:

# then unregister program file
y:vFNAME | e*:fUnRegisterProgram | v:vUnRegStatus

r@:{DELETED PROGRAM FILE: #vFNAME#
With Status: #vUnRegStatus#} | q!:


l:lNFINE
i!:{Invalid [blank] File Name!
PROGRAM FILE NOT DELETED.}
"
Y:fDeleteProgramFile
V@:dfDeleteProgramFile
L:lSYS_PROG_END
L:lSYS_SHELL_START
E*:fLogin
L:lSTART_Main
E*:fMakePS 
V@:dvLastPS
E*:fPROMPT
E*:fLogAI 
V:vIN_raw
t!.:
v:vIN_clean
Z:
v:vIN_clean_lc
y:vIN_clean_lc
F!:^create\s+.*$:lN_CRT_FILE
Y:vIN_raw
t!.:
V@:dvLastCommand
d:^\S+\s+
V:vFileName
E*:fCreateFile 
E*:fLogO
J:lFIN_CMD_PROCESSING
L:lN_CRT_FILE
F!:{^(ls|list|dir|files)$}:lN_LIST_FILES
V@:dvLastCommand
E*:fShowFilesList 
E*:fLogO
J:lFIN_CMD_PROCESSING
L:lN_LIST_FILES
F!:{^(ls|list|dir|files|find)\s+(\S+\s*)+$}:lN_SEARCH_FILES
Y:vIN_raw
t!.:
V@:dvLastCommand
d:^\S+\s+
V:vFilePattern
E*:fSearchFilesList 
E*:fLogO
J:lFIN_CMD_PROCESSING
L:lN_SEARCH_FILES
F!:{^(del|delete|rm)\s+(\S+\s*)+$}:lN_DEL_FILES
Y:vIN_raw
t!.:
V@:dvLastCommand
d:^\S+\s+
V:vFilePattern
E*:fDeleteFilesList 
E*:fLogO
J:lFIN_CMD_PROCESSING
L:lN_DEL_FILES
F!:^init\s+.*$:lN_INIT_FILE
Y:vIN_raw
t!.:
V@:dvLastCommand
d:^\S+\s+
V:vFileName
E*:fInitFile 
E*:fLogO
J:lFIN_CMD_PROCESSING
L:lN_INIT_FILE
F!:^write\s+.*$:lN_WRITE_FILE
Y:vIN_raw
t!.:
V@:dvLastCommand
d:^\S+\s+
V:vFileName
E*:fWriteFile 
E*:fLogO
J:lFIN_CMD_PROCESSING
L:lN_WRITE_FILE
F!:{^(show|cat|read)\s+.*$}:lN_READ_FILE
Y:vIN_raw
t!.:
V@:dvLastCommand
d:^\S+\s+
V:vFileName
E*:fReadFile 
E*:fLogO
J:lFIN_CMD_PROCESSING
L:lN_READ_FILE
F!:{^(exists)\s+.*$}:lN_EXISTS_FILE
Y:vIN_raw
t!.:
V@:dvLastCommand
d:^\S+\s+
V:vFileName
E*:fFileExists 
E*:fLogO
J:lFIN_CMD_PROCESSING
L:lN_EXISTS_FILE
F!:{^(rename|rn|mv|move)\s+([a-z0-9_]+(\.[a-z0-9]+)?)\s+([a-z0-9_]+(\.[a-z0-9]+)?)$}:lN_RENAME_FILE
Y:vIN_raw
t!.:
V@:dvLastCommand
d:^\S+\s+
V:vFileNames
Y:vFileNames 
D:\s+.*$ 
V:vFileName
Y:vFileNames 
D!:\s+.*$ 
V:vNewName
Y:vFileName 
E*:fFileExists
F!:{^FILE EXISTS$}:lCANNOT_RENAME
Y:vFileName 
E*:fProgramExists 
F:{^PROGRAM EXISTS$}:lPROCESS_RENAME_PROGRAM
Y:vFileName 
E*:fReadFile 
v@:dvDATABUFFER
Y:vNewName 
E*:fSaveDataToFile 
v:vWRITELOG
Y:vFileName 
E*:fDeleteFile
R@:{File #vFileName# RENAMED #vNewName#
With status: #vWRITELOG#}
E*:fLogO
J:lFIN_CMD_PROCESSING
L:lPROCESS_RENAME_PROGRAM
Y:vFileName 
E*:fReadFile 
v@:dvDATABUFFER
Y:vNewName 
E*:fSaveDataToFile 
v:vWRITELOG
Y:vNewName 
E*:fRegisterProgramFile 
v:vREGLOG
Y:vFileName 
E*:fDeleteProgramFile 
v:vDELLOG
R@:{Program File #vFileName# RENAMED #vNewName#
With [WRITE] status: #vWRITELOG#
And [REGISTER] status: #vREGLOG#
And [DEL] status: #vDELLOG#}
E*:fLogO
J:lFIN_CMD_PROCESSING
L:lCANNOT_RENAME
R@:{File #vFileName# does NOT EXIST!
RENAME Operation ABANDONED}
E*:fLogO
J:lFIN_CMD_PROCESSING
L:lN_RENAME_FILE
F!:{^(copy|clone)\s+([a-z0-9_]+(\.[a-z0-9]+)?)\s+([a-z0-9_]+(\.[a-z0-9]+)?)$}:lN_CLONE_FILE
Y:vIN_raw
t!.:
V@:dvLastCommand
d:^\S+\s+
V:vFileNames
Y:vFileNames 
D:\s+.*$ 
V:vFileName
Y:vFileNames 
D!:\s+.*$ 
V:vFileName2
Y:vFileName 
E*:fFileExists
F!:{^FILE EXISTS$}:lCANNOT_CLONE
Y:vFileName 
E*:fReadFile 
v@:dvDATABUFFER
Y:vFileName2 
E*:fSaveDataToFile 
v:vWRITELOG
R@:{File #vFileName# CLONED TO #vFileName2#
With status: #vWRITELOG#}
E*:fLogO
J:lFIN_CMD_PROCESSING
l:lCANNOT_CLONE
R@:{File #vFileName# does NOT EXIST!
CLONE Operation ABANDONED}
E*:fLogO
J:lFIN_CMD_PROCESSING
L:lN_CLONE_FILE
F!:^export\s+.*$:lN_EXPORT_FILE
Y:vIN_raw
t!.:
V@:dvLastCommand
d:^\S+\s+
V:vFileName
E*:fReadFile 
v:vLOG 
E*:fResetLog 
E*:fOSHandle 
X:{Thanks for Using }
V:vMSG
G*!:vDELIM2:vFileName:vLOG:vMSG
J:lEXIT_Main
L:lN_EXPORT_FILE
F!:^uimini$:lN_UIMINI
V@:dvLastCommand
V@:dvUIMODE_MINI:{1}
I!:{TOS Mini-Interface Mode Activated} 
E*:fLogO
J:lFIN_CMD_PROCESSING
L:lN_UIMINI
F!:^uilog$:lN_UILOG
V@:dvLastCommand
V@:dvUIMODE_MINI:{0}
I!:{TOS Log-Interface Mode Activated} 
E*:fLogO
J:lFIN_CMD_PROCESSING
L:lN_UILOG
F!:^about$:lN_ABOUT
V@:dvLastCommand
E*:fGetAboutTOS 
E*:fLogO
J:lFIN_CMD_PROCESSING
L:lN_ABOUT
F!:{^(help|man|whatis|info)$}:lN_HELP
V@:dvLastCommand
E*:fGetHelpMenu 
E*:fLogO
J:lFIN_CMD_PROCESSING
L:lN_HELP
F!:{^(help|man|whatis|info|about)\s+(\S+\s*)+$}:lN_HELP_TOPIC
Y:vIN_raw
t!.:
V@:dvLastCommand
d:{^\S+\s+} 
E*:fGetHelpOn 
E*:fLogO
J:lFIN_CMD_PROCESSING
L:lN_HELP_TOPIC
F!:^help-file[s]?$:lN_HELP_FILE
V@:dvLastCommand
E*:fGetFileCMDMenu 
E*:fLogO
J:lFIN_CMD_PROCESSING
L:lN_HELP_FILE
F!:{^help-prog([s]|rams)?$}:lN_HELP_PROG
V@:dvLastCommand
E*:fGetProgrammingMenu 
E*:fLogO
J:lFIN_CMD_PROCESSING
L:lN_HELP_PROG
F!:^help-util[s]?$:lN_HELP_UT
V@:dvLastCommand
E*:fGetUtilityMenu 
E*:fLogO
J:lFIN_CMD_PROCESSING
L:lN_HELP_UT
F!:^help-ui[s]?$:lN_HELP_UI
V@:dvLastCommand
E*:fGetUIMenu 
E*:fLogO
J:lFIN_CMD_PROCESSING
L:lN_HELP_UI
F!:^help-help[s]?$:lN_HELP_HELP
V@:dvLastCommand
E*:fGetHelpMenu 
X:{
To use or access help on any topic in the system, use any of the commands presented in the following Menu.
}
X!:{
Each of those commands leads you to further information about help on other commands. For example

> help help-util

Will explain what the help-util command is all about.}
E*:fLogO
J:lFIN_CMD_PROCESSING
L:lN_HELP_HELP
F!:^echo.*$:lN_ECHO
Y:vIN_raw
t!.:
d:^\S+\s*
V:vMESSAGE
X:{echo }
V@:dvLastCommand
Y:vMESSAGE 
E*:fLogO
J:lFIN_CMD_PROCESSING
L:lN_ECHO
F!:{^(fetch|web)\s+\S*$}:lN_FETCH
Y:vIN_raw
t!.:
V@:dvLastCommand
d:{^\S+\s+}
E*:fFixURL
W: 
E*:fLogO
J:lFIN_CMD_PROCESSING
L:lN_FETCH
F!:^tea\s+[\S]+.*$:lN_TEA
Y:vIN_raw
t!.:
d:^\S+\s+
V:vTEACode
X:{tea }
V@:dvLastCommand
Y:vTEACode
E: 
E*:fLogO
J:lFIN_CMD_PROCESSING
L:lN_TEA
F!:^math\s+[\S]+.*$:lN_MATH
Y:vIN_raw
t!.:
d:^\S+\s+
V:vMathExp
X:{math }
V@:dvLastCommand
Y:vMathExp
R.: 
E*:fLogO
J:lFIN_CMD_PROCESSING
L:lN_MATH
F!:^username\s+[\S]+.*$:lN_SETUSER
V@:dvLastCommand
Y:vIN_raw
t!.:
d:^\S+\s+
V@:dvUSERNAME 
E*:fLogO
J:lFIN_CMD_PROCESSING
L:lN_SETUSER
F!:{^(whoami|name)$}:lN_WHOAMI
V@:dvLastCommand
Y@:dvUSERNAME
E*:fLogO
J:lFIN_CMD_PROCESSING
L:lN_WHOAMI
F!:^time$:lN_TIME
v@:dvLastCommand
Z.:TIME 
E*:fLogO
J:lFIN_CMD_PROCESSING
L:lN_TIME
F!:^date$:lN_DATE
v@:dvLastCommand
Z.:DATE 
E*:fLogO
J:lFIN_CMD_PROCESSING
L:lN_DATE
F!:^timestamp$:lN_TSTAMP
v@:dvLastCommand
Z.:TIMESTAMP 
E*:fLogO
J:lFIN_CMD_PROCESSING
L:lN_TSTAMP
F!:{^(clear|reset)$}:lN_CLEAR
v@:dvLastCommand
E*:fResetLog
E*:fLogO
J:lFIN_CMD_PROCESSING
L:lN_CLEAR
F!:^print$:lN_PRINT
V@:dvLastCommand
E*:fGetLatestOutputBuffer
E*:fLogO
J:lFIN_CMD_PROCESSING
L:lN_PRINT
F!:^logout$:lN_LOGOUT
C!@:dvUSERNAME 
C*:vUSERNAME 
E*:fResetLog 
C@:dvLastOutput_backup
E*:fOSHandle 
X:{You have SIGNED OUT!
Thanks for Using } 
J:lEXIT_Main
L:lN_LOGOUT
F!:^exit$:lN_EXIT
E*:fResetLog 
E*:fOSHandle 
X:{Thanks for Using } 
J:lEXIT_Main
L:lN_EXIT
F!:^export$:lN_EXPORT
Y@:dvLOG
V:vLOG 
E*:fResetLog 
E*:fOSHandle 
X:{Thanks for Using }
V:vMSG
G*!:vDELIM2:vLOG:vMSG
J:lEXIT_Main
L:lN_EXPORT
F!:{^(runnable|register|add-program)\s+([a-z0-9_]+(\.[a-z0-9]+)?)$}:lN_RUNNABLE_FILE
Y:vIN_raw
t!.:
V@:dvLastCommand
d:^\S+\s+
V:vFileName
Y:vFileName 
E*:fFileExists
F!:{^FILE EXISTS$}:lCANNOT_REGISTER_PROG
Y:vFileName 
E*:fRegisterProgramFile
E*:fLogO
J:lFIN_CMD_PROCESSING
L:lCANNOT_REGISTER_PROG
R@:{File #vFileName# does NOT EXIST!
User-Program Registration Operation ABANDONED}
E*:fLogO
J:lFIN_CMD_PROCESSING
L:lN_RUNNABLE_FILE
F!:{^programs$}:lN_LIST_PROGRAMS
V@:dvLastCommand
E*:fShowProgramsList 
E*:fLogO
J:lFIN_CMD_PROCESSING
L:lN_LIST_PROGRAMS
F!:{^programs\s+(\S+\s*)+$}:lN_SEARCH_PROGRAMS
Y:vIN_raw
t!.:
V@:dvLastCommand
d:^\S+\s+
V:vFilePattern
E*:fSearchProgramsList 
E*:fLogO
J:lFIN_CMD_PROCESSING
L:lN_SEARCH_PROGRAMS
F!:{^(run|exec|execute)\s+.*$}:lN_RUN_FILE
Y:vIN_raw
t!.:
V@:dvLastCommand
d:^\S+\s+
V:vFileName
E*:fRunFile 
E*:fLogO
J:lFIN_CMD_PROCESSING
L:lN_RUN_FILE
F!:{^(isrunnable|canrun)\s+.*$}:lN_CANRUN_FILE
Y:vIN_raw
t!.:
V@:dvLastCommand
d:^\S+\s+
V:vFileName
E*:fProgramExists 
E*:fLogO
J:lFIN_CMD_PROCESSING
L:lN_CANRUN_FILE
F!:{^(del-prog|delete-program|unregister)\s+([a-z0-9_]+(\.[a-z0-9]+)?)$}:lN_UNREGISTER_FILE
Y:vIN_raw
t!.:
V@:dvLastCommand
d:^\S+\s+
V:vFileName
Y:vFileName 
E*:fFileExists
F!:{^FILE EXISTS$}:lCANNOT_REGISTER_PROG
Y:vFileName 
E*:fUnRegisterProgram
E*:fLogO
J:lFIN_CMD_PROCESSING
L:lCANNOT_REGISTER_PROG
R@:{File #vFileName# does NOT EXIST!
User-Program Unregistration Operation ABANDONED}
E*:fLogO
J:lFIN_CMD_PROCESSING
L:lN_UNREGISTER_FILE
F!:{^(update|update-program)\s+([a-z0-9_]+(\.[a-z0-9]+)?)\s+([a-z0-9_]+(\.[a-z0-9]+)?)$}:lN_UPDATE_PROGRAM
Y:vIN_raw
t!.:
V@:dvLastCommand
d:^\S+\s+
V:vFileNames
Y:vFileNames 
D:\s+.*$ 
V:vFileName
Y:vFileNames 
D!:\s+.*$ 
V:vUpdateFileName
Y:vUpdateFileName 
E*:fFileExists
F!:{^FILE EXISTS$}:lCANNOT_READ_SRC
Y:vFileName 
E*:fFileExists
F!:{^FILE EXISTS$}:lCANNOT_WRITE_TARGET
Y:vFileName 
E*:fProgramExists 
F:{^PROGRAM EXISTS$}:lPROCESS_UPDATE_PROGRAM
R@:{Target File #vFileName# IS NOT a PROGRAM!
Program UPDATE Operation ABANDONED}
E*:fLogO
J:lFIN_CMD_PROCESSING
L:lPROCESS_UPDATE_PROGRAM
Y:vUpdateFileName 
E*:fReadFile 
v@:dvDATABUFFER
Y:vFileName  
E*:fSaveDataToFile 
v:vWRITELOG
R@:{Program File #vFileName# UPDATED USING #vUpdateFileName#
With [WRITE] status: #vWRITELOG#}
E*:fLogO
J:lFIN_CMD_PROCESSING
L:lCANNOT_WRITE_TARGET
R@:{Target File #vFileName# does NOT EXIST!
Program UPDATE Operation ABANDONED}
E*:fLogO
J:lFIN_CMD_PROCESSING
L:lCANNOT_READ_SRC
R@:{Source File #vUpdateFileName# does NOT EXIST!
Program UPDATE Operation ABANDONED}
E*:fLogO
J:lFIN_CMD_PROCESSING
L:lN_UPDATE_PROGRAM
F!:{^([a-z0-9_]+(\.[a-z0-9]+)?)\s*$}:lN_RUN_FILE_PROG
Y:vIN_raw
t!.:
V@:dvLastCommand 
V:vFileName
Y:vFileName 
E*:fProgramExists 
F:{^PROGRAM EXISTS$}:lPROCESS_RUN_PROGRAM:lCMD_CATCHALL
L:lPROCESS_RUN_PROGRAM
Y:vFileName 
E*:fRunFile 
E*:fLogO
J:lFIN_CMD_PROCESSING
L:lN_RUN_FILE_PROG
L:lCMD_CATCHALL
F:^$:lFIN_CMD_PROCESSING
Y:vIN_raw 
V@:dvLastCommand 
I!:{--COMMAND NOT FOUND | INVALID COMMAND--}
E*:fLogO
J:lFIN_CMD_PROCESSING
L:lFIN_CMD_PROCESSING
E*:fLogAI
J:lSTART_Main
L:lEXIT_Main 
C!@:dvLOG 
Q!: 
L:lSYS_SHELL_END
   \end{lstlisting}
  \end{tcolorbox}
  \caption{TEA PROGRAM: TEA Operating System v1.1.0 (SANITIZED SOURCE)}
  \label{FIGTOSRI}
}